\chapter{Heat Cramps}

\section{Definition}
Heat cramps are painful, involuntary muscle spasms that occur during or after intense exercise in hot conditions. They are caused by electrolyte depletion, particularly sodium, from heavy sweating.

\section{Pathophysiology}
Prolonged sweating during physical activity in heat leads to loss of sodium and water. When sweat losses are replaced with only water, dilutional hyponatremia can develop, disrupting neuromuscular junction function and causing involuntary muscle contractions. The large muscles used during exercise are most commonly affected.

\section{Etiology}
\begin{itemize}
\item Intense physical activity in hot, humid conditions
\item Heavy sweating with inadequate salt replacement
\item Drinking only water without electrolytes during prolonged exercise
\item Sodium depletion
\item Poor physical conditioning
\item Inadequate acclimatization to heat
\item Fatigue
\item Dehydration
\end{itemize}

\section{Indications for Evaluation}
\begin{itemize}
\item Sudden onset of painful muscle cramps
\item Most common in calves, quadriceps, hamstrings, and abdominal muscles
\item Cramps during or after exercise
\item Heavy sweating prior to cramp onset
\item Muscles feel hard or knotted
\item Normal body temperature (differentiating from heat exhaustion)
\end{itemize}

\section{Related Symptoms}
\begin{itemize}
\item Heat exhaustion
\item Heat stroke
\item Dehydration
\item Muscle cramps (from other causes)
\item Exercise-associated muscle cramping
\end{itemize}

\section{Self-Care/Home Management}
\begin{itemize}
\item Stop activity and rest in a cool environment
\item Drink electrolyte-containing fluids (sports drinks, salt water)
\item Gently stretch and massage the cramped muscle
\item Apply firm pressure to the cramped muscle
\item Eat salty foods if able
\item Replace fluids containing electrolytes
\item Resume activity slowly and cautiously
\end{itemize}

\section{Diagnostic Approach}
\begin{itemize}
\item Clinical diagnosis based on history and physical exam
\item Electrolyte panel (may show hyponatremia or hypomagnesemia)
\item Creatine kinase if muscle injury suspected
\end{itemize}

\section{Treatment/Management Overview}
\begin{itemize}
\item Rest, stretching, and massage of the affected muscle
\item Oral electrolyte replacement
\item IV normal saline for severe or persistent cases
\item Prevention: proper hydration with electrolyte drinks
\item Acclimatize slowly to hot working conditions
\end{itemize}
